\documentclass[11pt]{article}
\usepackage[a4paper,margin=1in]{geometry}
\usepackage{fontspec}
\usepackage[boldfont=false]{xeCJK}
\setCJKmainfont{FandolSong-Regular.otf}
\usepackage{amsmath}
\usepackage{amssymb}
\usepackage{array}
\usepackage{booktabs}
\usepackage{graphicx}
\usepackage{adjustbox}
\usepackage{enumitem}
\setlist[itemize]{topsep=2pt,partopsep=0pt,parsep=0pt,itemsep=2pt,leftmargin=1.4em}
\setlist[enumerate]{topsep=2pt,partopsep=0pt,parsep=0pt,itemsep=2pt,leftmargin=1.6em}
\usepackage{parskip}
\usepackage{fvextra}
\fvset{breaklines=true,breakanywhere=true,fontsize=\small}
\setlength{\parindent}{0pt}

\begin{document}

\begin{center}
  {\LARGE\bfseries Deformation of solids}\\[0.35em]
  {\large A Level Physics}
\end{center}
\vspace{0.6em}

\section*{Forces that cause deformation}

When a force acts on a solid along its length, the object changes shape (\textbf{deformation} 形变). Two cases (treated as one-dimensional here):

\begin{itemize}
\item a \textbf{tensile} 拉伸 force stretches the object — it makes an \textbf{extension} 伸长量 $x$,

\item a \textbf{compressive} force squeezes the object — it makes a \textbf{compression} 压缩, treated as a negative extension.

\end{itemize}

The applied force is the \textbf{load} 负载. The change from the natural length is the extension (or compression).

\section*{Hooke's law and the spring constant}

For many materials at small extensions, \textbf{the extension is proportional to the load} — this is \textbf{Hooke's law} 胡克定律. The constant that links them is the \textbf{spring constant} 劲度系数 $k$:

\[
F = kx \qquad\Longleftrightarrow\qquad k = \frac{F}{x}.
\]

Unit of $k$: $\text{N m}^{-1}$.

\par\medskip
\noindent\begin{minipage}{\linewidth}
\centering
\adjustbox{max width=\linewidth, max totalheight=0.42\textheight}{%
  \includegraphics{/Users/joschka/Desktop/Programming/eXercise/output/eXam/handouts/a_level_physics/assets/06-p116-load-extension-graph.pdf}}\\[0.3em]
\emph{A load–extension graph: straight up to the limit of proportionality P, then it curves}
\end{minipage}\par\medskip

Reading a graph:

\begin{itemize}
\item A force–extension ($F$ against $x$) graph has gradient $k$ in the Hooke's-law region.

\item An extension–force ($x$ against $F$) graph has gradient $1/k$ in the Hooke's-law region.

\end{itemize}

A common trap: if a graph plots \textbf{length} $L$ against force, you can still find the spring constant from the gradient (since $L = L_{0} + F/k$, the gradient is $1/k$ — read it off carefully).

\subsection*{Limit of proportionality}

Hooke's law only holds up to the \textbf{limit of proportionality} 比例极限. Past this point the $F$ against $x$ line curves and is no longer straight. The material may still be \textbf{elastic} 弹性 (it returns to its first length when you remove the load) a little further, then it becomes \textbf{plastic} 塑性.

\subsection*{Springs in series and parallel}

You may need to combine spring constants:

\begin{itemize}
\item \textbf{Series} 串联 (one spring hangs from another): the same load passes through both, the total extension is the sum, so $\dfrac{1}{k_{\text{total}}} = \dfrac{1}{k_{1}} + \dfrac{1}{k_{2}}$.

\item \textbf{Parallel} 并联 (two springs side by side holding the same load): each takes half the load (if they are identical), the extensions are equal, so $k_{\text{total}} = k_{1} + k_{2}$.

\end{itemize}

\section*{Stress, strain and the Young modulus}

For a wire of uniform cross-section under a tensile load:

\begin{itemize}
\item \textbf{Stress} 应力 $\sigma = \dfrac{F}{A}$, where $F$ is the load and $A$ is the \textbf{cross-sectional area} 横截面积. Unit: $\text{Pa}$.

\item \textbf{Strain} 应变 $\varepsilon = \dfrac{x}{L_{0}}$, where $x$ is the extension and $L_{0}$ is the original length. Strain has no unit (it is a ratio of lengths).

\end{itemize}

The \textbf{Young modulus} 杨氏模量 is the ratio of stress to strain in the Hooke's-law region:

\[
E = \frac{\sigma}{\varepsilon} = \frac{F / A}{x / L_{0}} = \frac{F L_{0}}{A x}.
\]

Unit: $\text{Pa}$ (about $10^{11}$ for metals; e.g. steel $\approx 2.0 \times 10^{11}$ Pa).

\par\medskip
\noindent\begin{minipage}{\linewidth}
\centering
\adjustbox{max width=\linewidth, max totalheight=0.42\textheight}{%
  \includegraphics{/Users/joschka/Desktop/Programming/eXercise/output/eXam/handouts/a_level_physics/assets/06-p118-stress-strain-graph.pdf}}\\[0.3em]
\emph{A stress–strain graph, straight up to the limit of proportionality P}
\end{minipage}\par\medskip

The Young modulus is a property of the \textbf{material} — it does not depend on the wire's shape or size. The spring constant $k$ depends on both the material and the size: $k = EA/L_{0}$.

\subsection*{Experiment to find the Young modulus of a metal wire}

A standard setup:

\begin{enumerate}
\item Clamp one end of a long, thin wire to a fixed support. Pass the wire over a \textbf{pulley} 滑轮 at the edge of the bench so it hangs straight down.

\item Measure the original length $L_{0}$ between the clamp and a marker near the pulley, using a metre rule.

\item Measure the \textbf{diameter} 直径 $d$ of the wire at several places with a \textbf{micrometer} 螺旋测微器 and take the average. Work out $A = \pi d^{2}/4$.

\item Hang weights one at a time. Record the load $F$ and the extension $x$ (how far the marker moves against a fixed scale).

\item Plot $F$ against $x$. In the straight region the gradient is $EA/L_{0}$, so $E = \text{gradient} \times L_{0}/A$.

\end{enumerate}

Why a \textbf{long, thin} wire? To make the extension big enough to measure well. Why repeat readings and measure $d$ at several places? To reduce \textbf{random error} 随机误差 and check the wire is uniform.

\par\medskip
\noindent\begin{minipage}{\linewidth}
\centering
\adjustbox{max width=\linewidth, max totalheight=0.42\textheight}{%
  \includegraphics{/Users/joschka/Desktop/Programming/eXercise/output/eXam/handouts/a_level_physics/assets/06-p119-young-modulus-experiment.pdf}}\\[0.3em]
\emph{Apparatus for measuring the Young modulus of a wire}
\end{minipage}\par\medskip

\section*{Elastic and plastic behaviour}

As the load grows:

\begin{enumerate}
\item \textbf{Elastic and straight (Hooke obeyed)} — up to the limit of proportionality. Removing the load returns the object to its first length.

\item \textbf{Elastic but curved} — between the limit of proportionality and the \textbf{elastic limit} 弹性极限. The extension is no longer straight in the load, but on unloading the object still returns to its first length.

\item \textbf{Plastic} — past the elastic limit. On unloading, the object does \textbf{not} return to its first length; a permanent extension stays.

\end{enumerate}

Hooke's law only holds in the straight, elastic region.

\par\medskip
\noindent\begin{minipage}{\linewidth}
\centering
\adjustbox{max width=\linewidth, max totalheight=0.42\textheight}{%
  \includegraphics{/Users/joschka/Desktop/Programming/eXercise/output/eXam/handouts/a_level_physics/assets/06-p120-force-extension-elastic-plastic.pdf}}\\[0.3em]
\emph{Force–extension past the elastic limit: P and E marked, with a permanent extension B left after unloading}
\end{minipage}\par\medskip

On a force–extension graph for a material taken into the plastic region and then unloaded, the loading line and the unloading line are different. The unloading line is parallel to the first Hooke line but shifted to the right (the permanent extension left when the load reaches zero). The area between the loading and unloading lines is the energy turned into \textbf{thermal energy} 热能 in the material.

\section*{Energy stored in a stretched material}

The \textbf{work done} in stretching a material from $0$ to extension $x$, as the load grows from $0$ to $F$, is the \textbf{area under the force–extension graph}.

\par\medskip
\noindent\begin{minipage}{\linewidth}
\centering
\adjustbox{max width=\linewidth, max totalheight=0.42\textheight}{%
  \includegraphics{/Users/joschka/Desktop/Programming/eXercise/output/eXam/handouts/a_level_physics/assets/06-p121-energy-area-under-graph.pdf}}\\[0.3em]
\emph{The work done stretching a material is the area under the force–extension graph}
\end{minipage}\par\medskip

\subsection*{Hooke's-law material}

When Hooke's law holds, the $F$ against $x$ graph is a straight line through the origin. The area under it from $0$ to $x$ is a triangle:

\[
E_{\text{P}} = \tfrac{1}{2} F x = \tfrac{1}{2} k x^{2}.
\]

This is the \textbf{elastic potential energy} 弹性势能 stored in a spring or wire stretched within its limit of proportionality. An equal form:

\[
E_{\text{P}} = \frac{F^{2}}{2k}.
\]

\subsection*{Non-Hooke material}

For a graph that is not a straight line (a stretched rubber band, or a spring past its limit of proportionality), find the area by counting grid squares or by using \textbf{trapezia} 梯形. The same idea holds: \textbf{the area under the force–extension graph is the work done on the material}.

\subsection*{Comparing stored energy}

A common multiple-choice case: two materials are stretched by the same force, or by the same extension. Using $E_{\text{P}} = \tfrac{1}{2} F x$:

\begin{itemize}
\item same $F$, smaller $k$ (less stiff) → larger $x$ → more energy stored.

\item same $x$, larger $k$ (stiffer) → larger $F$ → more energy stored.

\end{itemize}

When a stretched spring is released onto a mass, the elastic potential energy becomes \textbf{kinetic energy} 动能 (and gravitational potential energy if the mass rises). Set $\tfrac{1}{2} k x^{2}$ equal to $\tfrac{1}{2} m v^{2}$ (plus any $mgh$) to find the speed or height.


\end{document}
